\documentclass[12pt]{article}
\usepackage[utf8]{inputenc}
\usepackage[T1]{fontenc}
\usepackage{hyperref}
\title{Tutong Language Resources for K-12 Curriculum Development}
\author{Compiled by Claude Sonnet 5 (claude-sonnet-5) for Andr\'as Kornai}
\date{}
\begin{document}
\maketitle
\section{Introduction}
Tutong (ISO 639-3 ttg) is an Austronesian language of the Rejang--Baram group spoken by roughly 17,000 people in Tutong District, Brunei. It is endangered, with shift toward Brunei Malay, Standard Malay and English.
\section{Primary Educational Resources}
No school primers or textbooks were found. A Tutong module has reportedly been taught at Universiti Brunei Darussalam since 2012 (per a secondary summary; not verified at a primary source).
\section{Linguistic and Reference Materials}
A Tutong--Malay, Malay--Tutong dictionary was reportedly published by Dewan Bahasa dan Pustaka in 1991; no online copy was found.
\begin{itemize}
\item Wikipedia, Tutong language: \url{https://en.wikipedia.org/wiki/Tutong_language}
\item Tutong Language Documentation Project (Research Data Australia): \url{https://researchdata.edu.au/tutong-language-documentation-project/12594}
\end{itemize}
\section{Religious and Community Resources}
None found in search results.
\section{Digital and Audio Resources}
The Tutong Language Documentation Project (PARADISEC collection BN1) holds audio and video recorded 2007--2009, mostly from speakers over 60, with dictionary elicitation items.
\begin{itemize}
\item PARADISEC item BN1-042: \url{https://catalog.paradisec.org.au/repository/BN1/042}
\item PARADISEC item BN1-027: \url{https://catalog.paradisec.org.au/repository/BN1/027}
\item PARADISEC item BN1-063: \url{https://catalog.paradisec.org.au/repository/BN1/063}
\end{itemize}
\section{Curriculum Development Context}
Tutong is not known to be taught in schools. The PARADISEC recordings and the dictionary are the main starting points, but no standardised orthography or pedagogical material was found.
\end{document}